\chapter{Conclusiones y trabajo futuro}
\label{ch:conclusiones}

\begin{chapterabstract}
Este capítulo cierra el proyecto contrastando lo realizado con el objetivo general, los objetivos específicos y el alcance definidos en el Capítulo~\ref{ch:introduccion}. Se concluye que el objetivo general se cumplió en el alcance de un MVP funcional que implementa la cadena completa de evidencia, desde el consentimiento del estudiante hasta la revisión docente, con un componente de IA no supervisado y sin veredictos automáticos, verificado mediante una validación técnica con datos sintéticos y pruebas automatizadas. Se distinguen los niveles de validación alcanzados, se sintetiza la viabilidad legal, técnica y económica del producto, factible bajo condiciones explícitas, se resumen los aportes del trabajo, se explicitan sus limitaciones y se ordena el trabajo futuro según su prioridad, comenzando por la medición del desempeño con datos de personas.
\end{chapterabstract}

\section{Cumplimiento del objetivo general}
\label{sec:concl_objetivo_general}

El objetivo general del proyecto era diseñar y desarrollar una plataforma de apoyo a la integridad académica en evaluaciones digitales, basada en la captura y el análisis de señales de interacción, para generar trazabilidad de autoría y alertas revisables. El objetivo se considera \textbf{cumplido} en el alcance de un producto mínimo viable. El diseño quedó documentado en los requerimientos (Capítulo~\ref{ch:requerimientos}) y en la arquitectura, el modelo de datos y el componente de IA (Capítulo~\ref{ch:tecnologias}); el desarrollo, en un MVP que funciona de punta a punta (Capítulo~\ref{ch:implementacion}): el docente configura qué se registra, el estudiante es informado y consiente, la plataforma captura solo las señales habilitadas, el \textit{backend} calcula métricas y genera indicadores trazables por reglas al finalizar cada sesión, al cerrar la evaluación compara cada sesión con la cohorte completa, y el docente revisa la evidencia y deja constancia de su revisión; los datos pueden eliminarse a pedido y se purgan al vencer el plazo de conservación. La validación técnica del Capítulo~\ref{ch:evaluacion} verificó que esa cadena produce el resultado esperado en 47 de 48 sesiones guionadas (12 de 12 sobre la versión final del código), ejecutadas en un navegador automatizado sobre el sistema real, que lo que muestra el panel docente coincide en todos los casos con lo que registra el sistema y que los eventos no contienen el texto de las respuestas. Una prueba adicional verificó que, ante cortes de red, la captura conserva lo registrado sin conexión o deja constancia de lo que pierde; las mediciones de rendimiento mostraron que se cumplen las cotas de procesamiento y de consulta, con la excepción de la entrega simultánea de 200 sesiones; y la evaluación experta y automatizada de usabilidad, con sus correcciones, llevó de 6 a 0 las tareas con problemas y de 15 a 0 las violaciones serias de accesibilidad (Tabla~\ref{tab:eval_usab_antes_despues}).

El MVP genera, así, las alertas revisables y trazables que el objetivo pedía: miden la atipicidad de la interacción respecto de la cohorte e indican al docente dónde mirar, sin probar la autoría de una respuesta ni la identidad de quien la escribe. La validación técnica se realizó con datos sintéticos generados por el equipo y escenarios guionados; medir \textbf{cuán bien} esas alertas distinguen conductas reales requiere datos de personas y es el primer punto del trabajo futuro (Sección~\ref{sec:concl_trabajo_futuro}).

\subsection{Niveles de validación}
\label{sec:concl_niveles}

El cumplimiento se ubica en los dos primeros de los tres niveles de la Tabla~\ref{tab:eval_niveles}. El \textbf{MVP} está implementado. La \textbf{validación técnica}, con datos sintéticos y pruebas automatizadas, está realizada y verifica el procedimiento y el funcionamiento interno. El \textbf{uso institucional} es el nivel siguiente y requiere un piloto con voluntarios, la calibración con datos reales, pruebas con docentes y estudiantes y las condiciones legales y técnicas de la Sección~\ref{sec:concl_viabilidad}.

\section{Cumplimiento de los objetivos específicos}
\label{sec:concl_objetivos_especificos}

La Tabla~\ref{tab:concl_objetivos} vincula cada objetivo específico con el capítulo donde se aborda, la evidencia que lo respalda y su grado de cumplimiento.

{\small
\begin{longtable}{>{\raggedright\arraybackslash}p{3.6cm} >{\raggedright\arraybackslash}p{7.4cm} >{\raggedright\arraybackslash}p{2.6cm}}
    \caption{Cumplimiento de los objetivos específicos}
    \label{tab:concl_objetivos} \\
    \toprule
    \tblhead \textbf{Objetivo específico} & \textbf{Evidencia} & \textbf{Cumplimiento} \\
    \midrule
    \endfirsthead
    \toprule
    \tblhead \textbf{Objetivo específico} & \textbf{Evidencia} & \textbf{Cumplimiento} \\
    \midrule
    \endhead
    \bottomrule
    \endfoot
    Relevar y sistematizar antecedentes académicos, técnicos y comerciales & Marco teórico (Sección~\ref{ch:marco_teorico}), estado del arte con tabla comparativa de enfoques y productos, y \textit{user research} con una encuesta a 213 estudiantes y tres entrevistas semiestructuradas (Capítulo~\ref{ch:user_research}). & Cumplido \\
    Determinar las señales de interacción y sus métricas & Correspondencia entre señales, métricas derivadas, información para el docente y método (Tabla~\ref{tab:senales_insights}); el MVP captura teclado, mouse, clics, \textit{scroll}, visibilidad, foco y portapapeles, y transforma en métricas la mayoría de ellas (Tabla~\ref{tab:impl_metricas}). & Cumplido en la identificación; implementado para la mayoría de las señales; pausas, tiempo por pregunta y tasa de corrección previstos como extensión \\
    Implementar y verificar un MVP que registre señales y señale sesiones atípicas & Captura embebida con envío por lotes e idempotencia, validación de eventos en el servidor, modelo de datos versionado, reglas por umbral al finalizar cada sesión e \textit{Isolation Forest} sobre la cohorte al cerrar la evaluación (Capítulo~\ref{ch:implementacion}); 178 pruebas automatizadas aprobadas, 142 del \textit{backend} y 36 del \textit{frontend}; evaluación reproducible con datos sintéticos, incluido el aporte del modelo sobre las reglas (nulo en los perfiles de pegado masivo, tecleo automatizado y cambio de pestaña; entre el 54 y el 67\,\% de detección en el perfil de ritmo distinto), prueba interna con 47 de 48 sesiones conformes, mediciones de rendimiento y evaluación experta y automatizada de usabilidad (Capítulo~\ref{ch:evaluacion}). & Cumplido en el nivel de validación técnica \\
    Presentar al docente indicadores explicables & Cada indicador de regla conserva la métrica, el valor y el umbral que lo originaron y se vincula a ellos en la base; el indicador del modelo se acompaña de las métricas con mayor desviación respecto de la cohorte, presentadas como dato descriptivo respecto de la cohorte; el panel docente los presenta en lenguaje claro, sin probabilidades de fraude (Secciones~\ref{sec:impl_pipeline} y~\ref{sec:impl_recorrido}); los indicadores se interpretan según las reglas del examen y se complementan con un panorama agregado de la evaluación (Sección~\ref{sec:impl_panorama}). & Cumplido para las reglas; para el modelo, con contexto descriptivo de la marca \\
    Incorporar información previa, consentimiento y política de captura desde el diseño & Política de captura versionada que define qué se registra, información previa y consentimiento obligatorio, eventos sin contenido escrito ni identificación de teclas, validados también en el servidor, revisión exclusivamente humana, eliminación de datos a pedido y purga al vencer el plazo de conservación (RF17, RNF10); análisis de viabilidad legal y descargo de responsabilidad (Capítulo~\ref{ch:viabilidad}). & Cumplido; el procedimiento institucional de derechos y la evaluación de impacto acompañan la puesta en producción \\
\end{longtable}
}

\section{Respeto del alcance}
\label{sec:concl_alcance}

El alcance excluía que ImprontIA funcionara como un sistema de detección automática de fraude, sancionara, bloqueara evaluaciones, reemplazara al docente o determinara si una respuesta fue producida por inteligencia artificial. El MVP respeta esas exclusiones por construcción: ningún componente emite un veredicto, los indicadores no tienen un campo de resultado, la captura registra acciones como el pegado o el cambio de pestaña sin impedirlas, y la única forma de registrar una conclusión es una revisión escrita por un docente, que no se puede editar después. La evaluación reforzó esta posición con datos sintéticos: el modelo marca alrededor del 10\,\% de las sesiones por diseño, lo que confirma que una marca de atipicidad funciona como una invitación a revisar.

\section{Evolución de la propuesta de valor}
\label{sec:concl_evolucion}

Los resultados de la evaluación permitieron precisar la propuesta de valor. La proporción de la cohorte que marca el modelo la fija aproximadamente un parámetro de diseño, con datos sintéticos el modelo concentró su aporte sobre las reglas en los desvíos de ritmo que ninguna regla observa, y la regla de ritmo uniforme mostró sensibilidad a una sola pulsación atípica (Capítulo~\ref{ch:evaluacion}): la atipicidad estadística resulta una guía para dirigir la revisión. En consecuencia, ImprontIA se presenta como una plataforma de \textbf{trazabilidad e inteligencia del proceso de evaluación digital}: hace visible cómo se produjeron las respuestas, interpreta los eventos según las reglas que el docente definió para cada examen y le ofrece un panorama agregado de la evaluación que resulta útil aunque nadie haya transgredido las reglas. El Incremento~12 incorporó esa ampliación sobre la misma arquitectura, el mismo cálculo de métricas y el mismo modelo, y con las reglas por defecto reproduce exactamente los resultados sintéticos (Sección~\ref{sec:impl_panorama}); su utilidad para los docentes se evaluará con ellos (Sección~\ref{sec:concl_trabajo_futuro}).

\section{Viabilidad del producto}
\label{sec:concl_viabilidad}

El análisis de viabilidad muestra que ImprontIA puede pasar de MVP a producto bajo condiciones explícitas y alcanzables. En lo legal, su uso es viable en Argentina si la institución actúa como responsable de los datos, informa y obtiene el consentimiento de los estudiantes y conserva la revisión humana como fundamento de toda decisión; el descargo de responsabilidad formaliza esa división (Capítulo~\ref{ch:viabilidad}). La eliminación de datos y la purga por plazo ya están implementadas; la evaluación de impacto en la protección de datos, el procedimiento para el ejercicio de los derechos de los titulares y la identificación del responsable en la vista del estudiante completan el cumplimiento. En lo técnico, la solución funciona con \textit{hardware} de uso general y sin infraestructura de aprendizaje automático dedicada, y los riesgos identificados son de madurez y no de factibilidad; las pruebas de carga en la infraestructura de producción, una auditoría de seguridad y la calibración con datos reales completan el paso a producción. En lo económico, el \refnegocio\ muestra que la rentabilidad depende de la adopción y del precio por asiento: con el precio de lista de USD~3,50 por asiento y una inversión inicial reducida, el escenario base es viable con un margen moderado (VAN de USD~6.201 y TIR de 23,8\,\% frente a una tasa exigida de 20,36\,\%), el optimista es claramente rentable y el pesimista no recupera la inversión; por eso se recomienda un lanzamiento y un financiamiento por etapas con un hito de continuidad (Sección~\ref{sec:neg_evaluacion}). Ninguna de las tres dimensiones presenta un impedimento de factibilidad. La legal y la técnica se completan con tareas acotadas de puesta en producción, cuyo costo principal ya contempla la inversión inicial del año~0, y la económica se confirma al validar el precio y la adopción con instituciones; ambos pasos forman parte del trabajo futuro.

\section{Aportes del trabajo}
\label{sec:concl_aportes}

\begin{itemize}
    \item \textbf{Un enfoque de trazabilidad del proceso, no de detección del producto.} Frente a los detectores de texto generado, cuyos sesgos y facilidad de evasión documenta la literatura~\parencite{Liang2023, Perkins2024}, y al \textit{proctoring} basado en cámara y bloqueo, ImprontIA analiza cómo se produce una respuesta con señales de baja invasividad y deja la interpretación al docente. El MVP muestra que ese enfoque es implementable con tecnologías de uso general.
    \item \textbf{Trazabilidad verificable de punta a punta.} Cada indicador puede seguirse hasta la métrica que lo originó y, por medio de ella, hasta los eventos de la sesión; la política aceptada por el estudiante queda fijada en su sesión. Es una respuesta concreta a la trazabilidad limitada que se observa en algunos sistemas de supervisión~\parencite{Heinrich2025}.
    \item \textbf{Una división de tareas entre reglas y modelo respaldada por la evaluación.} Las reglas explícitas cubren las señales discretas y el modelo poblacional cubre los desvíos simultáneos en varias métricas continuas; con datos sintéticos generados por el equipo y el protocolo vigente de cohorte al cierre, el modelo agrega detecciones exclusivamente en el perfil de ritmo distinto, que ninguna regla ve (del 54 al 67\,\%). Con el protocolo inductivo anterior, la combinación de ambos niveles dejó sin marcar el 3,6\,\% de las anomalías sintéticas.
    \item \textbf{Un procedimiento de evaluación reutilizable y reproducible.} El generador sintético, las métricas con convención explícita, el protocolo de repeticiones, las mediciones de rendimiento y la evaluación de usabilidad usan el mismo código que la plataforma y se reproducen con los comandos del Apéndice~\ref{apx:reproducibilidad}, por lo que el piloto con voluntarios podrá medirse sin construir herramientas nuevas.
    \item \textbf{Hallazgos técnicos documentados.} El mínimo de sesiones que necesita el modelo para separar atípicos, la dependencia del orden de entrega que introduce entrenarlo sesión por sesión, su insensibilidad a valores fuera del rango observado cuando se desvía una sola métrica, el hecho de que las desviaciones por métrica no explican sus marcas y la fragilidad de una regla basada en el desvío estándar son resultados observados que orientan cualquier implementación similar.
\end{itemize}

En términos de aporte a la sociedad, el proyecto propone una alternativa a la vigilancia invasiva para sostener la confianza en las evaluaciones digitales: menos datos, reglas visibles para el estudiante y decisiones siempre humanas. En el \textit{user research}, la aceptabilidad del registro de eventos de interacción fue heterogénea, mientras que el 93,0\,\% de los estudiantes encuestados asignó alta importancia a conocer previamente qué se registra y con qué finalidad y el 93,9\,\% a la existencia de revisión humana (Capítulo~\ref{ch:user_research}); el diseño incorpora esas garantías como requisitos y no como agregados.

\section{Limitaciones}
\label{sec:concl_limitaciones}

\begin{itemize}
    \item \textbf{Sin datos de personas.} La evaluación usó sesiones sintéticas generadas por el equipo y escenarios guionados. Valida el procedimiento y el funcionamiento interno, pero no estima el desempeño de detección con estudiantes, y los parámetros sintéticos no están calibrados con mediciones reales; varios perfiles se construyeron alrededor de los umbrales de las reglas.
    \item \textbf{Umbrales y corte del modelo provisorios.} Los umbrales de las reglas son marcadores de posición explícitos y la proporción de sesiones marcadas por el modelo es un parámetro elegido, no calibrado. La regla de ritmo uniforme, en particular, mostró ser sensible a una sola pulsación atípica.
    \item \textbf{Observación limitada a la plataforma.} La captura embebida solo observa la pestaña de ImprontIA, que por eso aloja la superficie de respuesta; la integración con otras plataformas de examen se plantea como trabajo futuro.
    \item \textbf{Aspectos no evaluados con personas.} La comprensión de los indicadores por parte de docentes y la interferencia percibida por los estudiantes no se midieron; la evaluación de usabilidad fue experta y automatizada, sin participantes. El comportamiento con muchas sesiones concurrentes se midió en una sola máquina, y la entrega simultánea de 200 sesiones superó la cota de 5\,s.
    \item \textbf{Infraestructura de producción.} El \textit{backend} se verificó en modo producción sobre PostgreSQL y el esquema se gestiona con migraciones versionadas, y la demostración funciona en la nube (Vercel y Neon, en planes gratuitos y con datos sintéticos; Sección~\ref{sec:impl_acceso}), pero el despliegue institucional sigue pendiente: no hay integración continua ni pruebas de carga en la infraestructura de producción, los contenedores Docker solo se verificaron en su configuración y no se realizó una auditoría de seguridad; el procedimiento institucional para ejercer los derechos de los titulares se define con cada institución.
    \item \textbf{Supuestos del modelo de negocio.} Los precios, la adopción y los costos del análisis económico-financiero se apoyan en referencias públicas de mercado y no en una validación con instituciones; la disposición a pagar no se relevó.
\end{itemize}

\section{Trabajo futuro}
\label{sec:concl_trabajo_futuro}

El trabajo futuro se ordena por prioridad, de modo que cada paso habilite el siguiente:

\begin{enumerate}
    \item \textbf{Realizar simulaciones controladas y el piloto con voluntarios}: primero, simulaciones del propio equipo con las herramientas del Incremento~15 (Sección~\ref{sec:impl_simulacion}), para ensayar el procedimiento con datos de interacción reales; después, el piloto con el protocolo de la Sección~\ref{sec:eval_piloto_futuro}: diseño intrasujeto con conductas indicadas por consigna, partición por participante e intervalos por remuestreo, reutilizando el código de evaluación existente. Las simulaciones están previstas para diciembre de 2026 y el piloto para el primer cuatrimestre de 2027, según el cronograma de la Sección~\ref{sec:eval_piloto_futuro}; incluye medir con docentes cuántas sesiones requieren revisión y cuánto tiempo insume cada una, la pregunta que la evaluación sintética deja abierta (Sección~\ref{sec:eval_ahorro_docente}).
    \item \textbf{Calibrar umbrales y modelo} con los datos del piloto: reemplazar el desvío estándar de la regla de ritmo por una medida robusta de dispersión, como la desviación absoluta mediana~\parencite{Leys2013}, excluir las pausas por pérdida de foco del tiempo entre teclas, evaluar la pérdida de foco de ventana como señal y ajustar los cortes provisorios: la proporción de sesiones que marca el modelo, el percentil que define el nivel alto y los parámetros de la sincronización.
    \item \textbf{Evaluar con docentes} la comprensión de los indicadores, el esfuerzo de revisión (RNF06, RNF08) y la utilidad del panorama de la evaluación y de las reglas del examen (RF15, RF16), con un instrumento estandarizado de usabilidad percibida como la escala SUS~\parencite{Bangor2008}, y con estudiantes la comprensión de la información previa y la interferencia percibida (RNF02, RNF07).
    \item \textbf{Validar el modelo de negocio} con instituciones: relevar la disposición a pagar con métodos que no dependan solo de preguntarla en forma directa, cuyas respuestas tienden a estar sesgadas~\parencite{Breidert2006}, y el precio por asiento aceptable, que es la variable que más influye en la rentabilidad, antes de comprometer la inversión prevista.
    \item \textbf{Completar las variables derivadas} previstas: tasa de corrección, velocidad de tipeo y ráfagas, pausas y aceleración del mouse, \textit{scroll} y tiempo por pregunta, con soporte de evaluaciones de varios ítems.
    \item \textbf{Preparar la operación en producción}: despliegue institucional a partir de la demostración ya publicada y del \textit{script} de verificación, integración continua, persistencia en el navegador de los eventos pendientes de envío, puntaje del modelo por lotes para acortar el cierre de evaluaciones grandes, ajuste del mantenimiento de las conexiones (\textit{keep-alive}) y de la entrega simultánea, pruebas de carga en la infraestructura de producción, auditoría de seguridad, procedimiento para los derechos de los titulares y evaluación de impacto en la protección de datos.
    \item \textbf{Explorar la integración} con plataformas de gestión del aprendizaje y, si se cuenta con datos etiquetados suficientes, modelos supervisados dentro del mismo ecosistema, como prevé la Sección~\ref{sec:procesamiento_ia}.
\end{enumerate}

\section{Reflexión final}

El problema que dio origen a este proyecto no es técnico en su raíz: es la pérdida de visibilidad sobre cómo se produce una respuesta académica cuando la evaluación se vuelve digital y las herramientas de generación de texto están al alcance de todos. En lugar de una respuesta automática, ImprontIA devuelve al docente información sobre el proceso, con reglas conocidas por el estudiante y con la decisión siempre en manos de una persona. El MVP demuestra que esa propuesta puede construirse y la validación técnica, que funciona como se diseñó; el piloto con voluntarios es el paso que permitirá medir cuánto aporta en la práctica.
